% !TeX root = ../Thesis.tex
\chapter{研究方案}
\label{chap:research-proposal}

\section{研究目标}

本研究依据《义务教育物理课程标准（2022年版）》科学思维与探究实践核心素养要求，针对初中物理典型迷思概念识别滞后、干预针对性不足、引导连续性弱等问题，以及生成式人工智能教育应用中存在的“过度代答”风险，以苏格拉底式对话教学法为基础，结合大语言模型与Agent编排技术，设计并实现面向初中物理典型迷思概念干预的苏格拉底式对话教育智能体，以电学与浮力为主要验证场景开展研究。具体研究目标如下：

\begin{enumerate}
    \item 构建面向初中物理典型迷思概念干预的苏格拉底式对话教学策略体系。围绕电学与浮力典型迷思概念，梳理学生错误认知的表现、成因与干预需求，提炼适用于概念澄清、假设挑战、证据追问、后果探索与支架支持的苏格拉底提问策略，形成适配初中物理教学场景的策略分类体系，探索其可计算化表达方式。

    \item 设计适配苏格拉底对话逻辑的教学状态机模型与教育智能体原型。基于LangGraph框架，将苏格拉底式教学“诊断—冲突—支架—验证”的动态引导逻辑转化为可执行的教学状态机模型，构建具备状态流转、上下文记忆、策略调用与知识检索能力的教育智能体原型，提升对话过程的长程逻辑一致性、教学可控性与干预针对性。

    \item 构建教育场景专用的防答案泄露护栏机制。针对初中教育场景中生成式人工智能可能出现的“过度帮助”“直接代答”风险，设计兼顾输入识别、生成约束、输出审核与重生成控制的多层护栏机制，确保智能体交互始终遵循“引导而非代答”原则，提升系统伦理安全性与教学规范性。

    \item 验证苏格拉底式对话教育智能体的有效性、安全性与教学适配性。通过仿真测试、专家评估与小规模试用，从迷思概念识别与修正效果、对话教学适切性、答案泄露防控能力、系统运行稳定性等维度，对教育智能体原型进行综合评估，检验研究方案的可行性。
\end{enumerate}

\section{研究内容}

围绕上述目标，本研究从教学策略建模、系统架构设计、安全机制构建与效果评估四个维度展开，形成“问题识别—策略建模—系统实现—效果验证”的研究路径。考虑硕士论文研究周期与开发评估成本，本研究以初中物理中迷思特征典型、教学边界清晰的电学与浮力主题为主要验证场景，研究结论优先适用于同类概念性干预任务。

\subsection{初中物理典型迷思概念及苏格拉底式对话教学策略体系构建}

本研究首先聚焦初中物理典型迷思概念，重点围绕电学与浮力主题，结合课程标准、教材内容、教学案例与相关研究文献，梳理学生常见错误认知的表现、形成特点及其与科学概念的偏差。在此基础上，结合苏格拉底式教学法核心思想，分析教师概念干预过程中的追问方式、支架方式与引导逻辑，归纳形成适用于初中物理典型迷思概念干预的苏格拉底式对话教学策略体系。

本研究将对上述策略进行结构化表达，明确不同策略的教学目标、适用情境、触发条件、话语功能及其与迷思概念类型的对应关系，为后续教学状态机设计与教育智能体实现提供规则基础。该部分核心任务是探索如何将苏格拉底式教学策略从经验性教学智慧转化为可建模、可调用、可执行的系统规则。

\subsection{基于LangGraph的教学状态机设计与教育智能体实现}

在策略体系构建基础上，本研究基于LangGraph框架设计适配苏格拉底对话逻辑的教学状态机模型。该模型围绕“迷思识别—澄清追问—假设挑战—证据追问—支架支持—理解验证”等核心教学环节，定义核心状态、状态转移条件、循环与回退机制，使智能体能够根据学生回答动态调整对话路径，而非采用无状态、线性的问答模式。

系统实现层面，本研究以大语言模型为语言理解与生成基础，以教学状态机为对话控制核心，整合初中物理领域知识库、提问策略库、对话日志记录与知识检索模块，实现面向初中物理典型迷思概念干预的苏格拉底式对话教育智能体原型。该原型不追求复杂的通用Agent功能，重点通过显式状态控制提升对话的教学可控性、逻辑一致性与干预针对性。

\subsection{教育场景专用的防答案泄露护栏机制设计}

针对生成式人工智能教育应用中易直接提供结论、削弱学生独立思考过程的问题，本研究构建面向教学场景的防答案泄露护栏机制。该机制遵循“教学支持优先、直接代答受限”原则，从输入识别、生成约束、输出审核与必要时重生成四个环节设计。

具体而言，输入环节识别学生是否存在直接索要答案、规避思考过程等高风险意图；生成环节通过规则约束与教学提示控制模型直接输出结论的倾向；输出环节对候选回复进行审核，识别是否存在过度提示、直接泄露答案或替代学生完成推理的风险，必要时触发重生成。通过上述多层约束，提升系统教育伦理安全性，确保其服务于“引导学习”而非“替代学习”。

\subsection{教育智能体有效性、安全性与教学适配性评估}

为检验研究设计方案的可行性，本研究构建多维评估框架，对教育智能体开展系统验证，评估围绕有效性、安全性与教学适配性三个核心维度推进。其中，有效性评估重点考察教育智能体能否准确识别学生迷思概念，通过多轮对话促进学生修正错误认知，形成更接近科学概念的解释；安全性评估重点考察系统在面对直接索要答案、诱导越界回答等情境时，能否稳定执行拒绝代答、转向引导的策略，控制答案泄露风险；教学适配性评估重点考察系统提问是否符合初中物理教学规律，支架提供是否适度，对话节奏是否符合学生认知特点，整体交互是否具备教育可接受性。评估方式综合采用模拟学生对话测试、专家审阅与小规模试用，结合基线对照与模块消融，形成对系统原型的全面验证依据。

\section{拟解决的关键问题}

围绕研究目标与内容，本研究重点解决以下四个关键问题：

\begin{enumerate}
    \item 苏格拉底式对话教学策略面向初中物理典型迷思概念干预的可计算化表达。
    苏格拉底式教学法是高度依赖教师教学经验与情境判断的对话教学方式，具有动态性与隐性特征。如何将其转化为适配初中物理典型迷思概念干预的结构化策略体系，并进一步表达为系统可识别、可调用的规则模型，是本研究需解决的首要问题。

    \item 兼顾教学逻辑一致性与动态适应性的对话控制机制构建。
    典型迷思概念干预通常无法通过单轮问答完成，需经历澄清、追问、挑战、支架与验证等多轮递进过程。传统对话系统易出现上下文遗忘、逻辑断裂或过早给出答案的问题。如何借助教学状态机实现对话过程的显式控制，使教育智能体在多轮交互中保持教学目标一致、引导路径清晰，并根据学生反馈动态调整，是本研究需解决的第二个关键问题。

    \item 教育场景中答案泄露防控与教学支持功能的平衡。
    对教育智能体而言，“不给答案”不等于“拒绝帮助”，核心是在安全约束与教学支持之间取得平衡。如何构建兼具识别、约束、审核与重生成能力的护栏机制，使系统既能避免直接代答，又能提供恰当的启发与支架，是本研究需突破的第三个关键问题。

    \item 兼顾有效性、安全性与教学适配性的评估框架建立。
    苏格拉底式对话教育智能体的价值不仅体现在技术可运行性上，更体现在是否真正促进学生概念理解、符合教学规范、具备伦理安全性。如何构建多维度、可操作的评估方案，通过仿真测试、专家评估、小规模试用与基线对照对系统原型进行科学验证，是本研究拟解决的第四个关键问题。
\end{enumerate}

\section{拟采取的研究方法}

本研究聚焦初中物理典型迷思概念（以电学与浮力为例）的苏格拉底式对话教育智能体设计与实现，研究内容涵盖教学策略建模、系统原型开发与效果验证，采用教育研究与设计科学研究结合的综合研究思路。整体以设计科学研究（Design Science Research, DSR）为总体范式，按照“问题识别—方案设计—原型实现—效果评估—迭代优化”路径推进，确保研究既回应真实教学问题，又形成可验证的技术人工物。

结合研究目标与任务，采用以下研究方法：


\paragraph{设计科学研究法。}

设计科学研究法是本研究的总体方法论框架。本研究所构建的苏格拉底式对话教育智能体，本质上是服务于教学问题解决的技术人工物，研究重点不仅在于系统能否运行，更在于能否有效回应初中物理典型迷思概念干预这一真实教育问题。因此本研究遵循设计科学研究基本逻辑：首先识别初中物理典型迷思概念干预中的核心教学问题，明确系统设计目标；其次围绕教学策略体系、教学状态机与防答案泄露机制进行方案设计；再次基于大语言模型与图编排框架实现系统原型；最后通过仿真测试、专家评估与小规模试用对系统原型进行验证，根据评估结果优化迭代。

\paragraph{内容分析法与案例分析法。}

内容分析法与案例分析法用于提炼适配初中物理典型迷思概念干预的苏格拉底式对话教学策略。研究选取课程标准、教材习题解析、优质课堂实录、典型教学案例以及已有研究中与概念干预相关的师生对话材料作为分析对象，围绕“澄清追问、挑战假设、证据追问、后果探索、类比支架、理解验证”等功能维度建立编码框架，对教学行为进行系统编码分析。通过归纳不同提问策略的适用情境、触发条件、教学意图与典型表达方式，逐步形成面向初中物理典型迷思概念干预的苏格拉底式对话教学策略体系，为后续教学状态机建模提供规则基础。

\paragraph{原型迭代开发法。}

原型迭代开发法用于完成教育智能体系统的设计、实现与优化。研究以“最小可行原型（MVP）—局部测试—问题修正—功能扩展”的方式推进系统开发，优先实现教学状态流转、提问策略调用与基本对话控制功能，再逐步补充领域知识检索、防答案泄露护栏、对话日志记录与评估模块。该方法有助于在开发过程中及时发现系统在逻辑流转、上下文保持、策略调用与输出控制等方面的问题，通过多轮迭代提升系统的稳定性、可控性与教学适配性。

\paragraph{仿真实验法。}

仿真实验法用于在可控条件下检验系统原型在不同迷思概念情境中的运行表现。研究构建具有不同认知水平、表达风格与错误概念特征的模拟学生，与教育智能体进行多轮自动化对话测试，重点考察系统对典型迷思概念的识别能力、引导过程的连续性、状态流转的合理性以及防答案泄露机制的稳定性。同时通过设置基线版本与模块消融版本，对显式状态控制和护栏机制的作用进行比较分析。

\paragraph{专家评估法。}

专家评估法用于考察系统提问策略的教学合理性、学科准确性与支架使用的适切性。研究邀请3—5位具有初中物理教学经验或科学教育研究背景的专家，对系统生成的典型对话案例进行审阅评价，重点关注提问是否符合初中生认知特点、能否有效暴露并干预迷思概念、是否体现“引导而非代答”原则，以及整体交互是否具备教育可接受性。专家评估结果作为检验系统教学适配性的核心依据。

\paragraph{小规模试用法。}

在系统原型经过前期测试与专家审阅后，本研究在伦理审批通过前提下开展小规模试用，观察系统在真实用户环境中的交互效果。试用对象为少量目标用户，重点关注学生使用过程中的理解程度、交互体验、接受意愿与学习支持感，通过试用反馈发现系统在语言表达、提问节奏、支架方式与交互流程中的不足。小规模试用为探索性验证，不用于推断大规模教学效果。


\section{技术路线}

本研究的技术路线遵循 “教学问题驱动、教育理论支撑、系统设计实现、实验验证反馈” 的核心逻辑，构建形成从理论建模到系统开发、再到效果评估的完整闭环研究路径，所设计的苏格拉底式对话教育智能体设计框架详见附录 \ref {Detailed_Design_Framework}，该路线具体分为以下四个阶段：

\paragraph{第一阶段：初中物理典型迷思概念与苏格拉底式教学策略建模。}

本阶段完成研究的理论建构与规则提炼。首先，结合课程标准、教材与相关研究成果，梳理初中物理典型迷思概念，重点聚焦电学与浮力主题，形成迷思概念条目库。其次，通过文献研究、案例分析与内容分析，提炼教师概念干预中常用的苏格拉底式提问方式，归纳其教学功能、触发条件与对话特征，构建苏格拉底式对话教学策略体系。最终将教学策略转化为可供系统调用的结构化规则，为后续教学状态机设计与系统实现提供依据。

\paragraph{第二阶段：教学状态机设计与教育智能体原型开发。}
本阶段聚焦系统核心架构的设计与实现。研究基于LangGraph框架，围绕“迷思识别—澄清追问—假设挑战—证据追问—支架支持—理解验证”等核心教学环节，设计教学状态机模型，明确状态定义、转移条件、循环与回退机制。原型实现层面，系统以大语言模型为语言生成基础，以教学状态机为对话控制核心，整合初中物理领域知识库、提问策略库、对话日志模块与知识检索模块，构建面向典型迷思概念干预的苏格拉底式对话教育智能体原型。

\paragraph{第三阶段：防答案泄露护栏机制开发与系统集成。}

本阶段重点解决教育场景中生成式人工智能“过度解答”与“直接代答”的风险问题。研究围绕“输入识别—生成约束—输出审核—必要时重生成”思路，设计教育场景专用的防答案泄露护栏机制。输入侧识别用户是否存在直接索要结论、规避思考过程等高风险意图；生成侧通过系统规则与提示约束限制直接输出答案；输出侧对候选回复进行审核，判断是否存在泄露标准答案、替代学生推理或超出教学引导边界的风险。完成上述机制设计后，将其与教学状态机、知识库、大语言模型集成为统一的系统原型。




\paragraph{第四阶段：实验评估与迭代优化。}
\label{subsec:evaluation_and_iteration}

本阶段对系统原型进行多维度测试与迭代优化。研究通过功能测试、仿真实验、专家评估与小规模试用收集系统在不同场景下的运行数据与反馈信息，围绕迷思概念识别效果、引导过程质量、答案泄露控制能力、对话逻辑一致性与教学适配性等维度进行分析。同时设置基线对照与模块消融，比较显式状态控制与护栏机制的实际作用。根据评估结果修正提问策略、状态转移逻辑与护栏规则，优化系统原型，形成稳定的研究成果。

\section{实验方案}

为验证所设计的苏格拉底式对话教育智能体对初中物理典型迷思概念的干预效果，本研究构建兼顾有效性、安全性与教学适配性的实验方案。实验遵循“分层验证、逐步推进”思路，先验证系统功能与安全边界，再验证教学合理性，最后进行初步应用测试。考虑研究边界与可操作性，实验以电学与浮力主题的典型迷思概念为核心场景。

\subsection{实验对象与实验场景}

本研究选择初中物理典型迷思概念作为实验场景，优先选取电学中的“电流会被灯泡消耗”“电池中的电会被用完”，以及浮力中的“浮力大小只由物体轻重决定”“浮力等于重力”等迷思概念。选择上述主题的原因在于：一是这类迷思概念在初中物理教学中普遍存在，二是其科学概念边界清晰，便于构建干预脚本与评估标准。

实验对象涵盖三类群体：用于大规模仿真对话测试的模拟学生、用于教学适配性评估的初中物理教学或科学教育研究领域专家、伦理审批通过后参与小规模试用的目标用户。

\subsection{实验材料与数据来源}

实验材料涵盖五大类：基于教材、课程标准与相关研究构建的初中物理典型迷思概念条目库，经内容分析形成的苏格拉底式提问策略库，用于系统知识支撑的初中物理领域知识片段，用于测试的模拟学生画像、测试问答脚本与对抗性输入样例，以及专家评估量表、小规模试用反馈表与对话日志记录模板。上述材料共同为系统开发、对话测试与评估分析提供支撑。

\subsection{实验步骤设计}

本研究实验分为四个层次推进：首先是功能与安全测试，系统原型开发完成后，通过预设测试样例验证状态机能否按预期完成状态流转、知识检索是否正确触发、护栏机制能否识别高风险输入并抑制直接代答行为，排查系统逻辑错误与运行缺陷。其次是仿真对话实验，通过构建具有不同认知水平、表达风格与迷思概念特征的模拟学生，与教育智能体进行多轮自动化对话，重点观察系统能否准确识别不同类型的迷思概念、通过多轮追问促使学生修正错误认知，以及面对直接索要答案等诱导性提问时能否稳定执行拒答并转向启发式引导，该环节可在低伦理风险、高可控性条件下对系统进行批量测试与压力测试。随后开展基线对照与模块消融实验，设置基线版本与模块消融版本进行比较：基线版本采用仅依赖提示词控制的普通对话版本，消融版本分别去除教学状态机或护栏机制，以观察显式状态控制与防答案泄露机制对系统表现的具体影响。最后进行专家评估与小规模试用：在仿真测试基础上，邀请初中物理教学或科学教育研究领域专家对系统对话案例进行审阅评分，重点考察提问策略的教学适切性、科学准确性、支架使用恰当性与整体交互的教育合理性；伦理审批通过且条件允许的前提下，组织小规模试用，由少量目标用户与系统进行真实交互，收集关于可理解性、可接受性、交互流畅性与学习支持感等方面的反馈，该阶段为探索性应用验证，不用于推断大规模教学效果。

\subsection{评估指标设计}

为全面评估系统表现，本研究从四个维度构建指标体系：一是有效性指标，包括迷思概念识别准确率、引导后认知修正率、单次有效干预平均对话轮数、系统能否推动学生生成更接近科学概念的解释等，其中迷思概念识别准确率通过系统识别结果与人工标注结果的一致程度衡量，认知修正率根据对话前后学生表述的变化判定。二是安全性指标，包括面对直接索要答案时的拒答成功率、护栏拦截率、答案泄露率、越界回复发生率等，用于评估系统的教育安全边界控制能力，其中答案泄露率主要考察输出中是否出现直接结论、完整答案或替代学生完成推理的内容。三是教学适配性指标，包括提问是否符合初中生认知特点、支架提供是否适度、追问节奏是否合理、是否体现“引导而非灌输”原则等，相关数据通过专家量表评分与典型案例分析获得。四是系统运行指标，包括状态流转成功率、上下文保持连续性、响应延迟是否在可接受范围、系统是否存在明显逻辑死循环等，用于评价系统原型的稳定性与可用性。

\subsection{数据分析方式}

实验数据综合采用定量与定性分析方式处理：对于仿真测试、基线对照、小规模试用中产生的量化指标，采用描述性统计方法汇总比较，重点呈现不同版本系统在有效性、安全性与稳定性方面的差异；对于专家评语、典型对话案例与系统异常日志，采用案例分析与内容分析方法归纳，揭示系统在提问策略执行、状态流转控制与护栏约束方面的优势与不足。通过定量结果与定性分析结合，全面说明研究方案的实际效果。

\section{安全论证}

本研究从教育伦理安全、内容生成安全、未成年人数据隐私安全与研究过程风险控制四个维度构建安全保障方案。教育伦理安全层面，通过多层防答案泄露护栏机制约束智能体“引导而非代答”的核心行为，避免系统成为学生规避思考、直接获取作业答案的工具；内容生成安全层面，通过教学状态机与RAG检索增强机制的双重约束，将智能体的知识边界限定在初中物理课程标准与教材范围内，同时对输入与输出内容进行双向过滤，降低生成内容出现学科性错误与不当输出的风险；未成年人数据隐私安全层面，研究前期以模拟学生生成的合成数据为主要测试数据，如后期开展小规模试用，将遵循知情同意、匿名化处理、最小化数据采集与研究目的限定原则，不采集与研究无关的个人敏感信息；研究过程风险控制层面，通过预设样例的多轮内部测试、专家对系统输出的审查，以及小规模试用中的人工观察与必要干预机制，及时发现并修正系统边界问题与不当输出。

本研究严格遵循K12教育场景未成年人保护、知情同意与教育伦理要求开展研究；如涉及真实用户试用，将在学校或学院伦理审批程序通过后实施。

\section{可行性分析}

本研究在理论、技术、数据资源与实践应用四个维度具备可行性。理论层面，苏格拉底式教学法、概念转变理论、建构主义学习理论与支架式教学理论均有成熟的学术支撑，且初中物理电学与浮力模块的迷思概念已有丰富的研究积累，为本研究提供理论依据；技术层面，当前中文大语言模型已具备良好的语言理解与生成能力，LangGraph等图式编排框架能够支持对话状态流转与持久化管理，向量数据库与护栏工具也为系统开发提供成熟的技术条件；数据与资源层面，人教版初中物理教材、义务教育物理课程标准与相关学术文献均可作为知识库与策略建模的核心数据来源，依托培养单位的学科平台可便捷获取专家支持；实践层面，本研究聚焦初中物理典型迷思概念干预这一真实教学问题，研究对象明确、场景可控，具备应用探索价值。

本研究存在的潜在风险包括：长程对话中状态判断可能出现偏差、护栏机制可能误伤正常教学支持、小规模试用样本有限等。对此，研究将通过对话日志复核、规则迭代优化、专家审阅与探索性验证等方式控制，提升研究的稳健性。

\section{本研究的特色与创新之处}

本研究的特色与创新点主要包括三方面：一是面向典型迷思概念干预的教学策略可计算化表达，本研究将初中物理典型迷思概念干预中的苏格拉底式提问策略，从经验性教学智慧转化为可建模、可调用、可执行的结构化策略体系，为教育智能体的对话逻辑提供教学理论支撑的规则基础。二是基于显式教学状态机的对话控制机制，本研究围绕“迷思识别—澄清追问—假设挑战—证据追问—支架支持—理解验证”的教学逻辑构建显式状态机，增强多轮对话中的逻辑一致性、过程可控性与干预针对性，降低传统无状态对话系统中上下文遗忘与逻辑断裂的问题。三是面向教育场景的防答案泄露护栏设计，本研究设计兼顾输入识别、生成约束、输出审核与必要时重生成的多层护栏机制，在避免系统直接代答的同时保留必要的启发与支架支持，实现教学支持与教育安全的平衡。

\section{预期的论文进展和成果}


\begin{table}[htbp]
    \centering
    \caption{预期进展时间表}
    \label{tab:progress-timetable}
    \begin{tabular}{l p{\dimexpr\textwidth - 2.5cm - 4\tabcolsep\relax}}
        \hline
        时间 & 预期主要进展内容 \\
        \hline
        2026年2月 & 完成文献调研与研究方案细化，分析《义务教育物理课程标准（2022年版）》与初中物理教材，构建以电学与浮力为例的“典型迷思概念条目库”和“苏格拉底式提问策略库”。 \\
        \hline
        2026年3月 & 完成基于LangGraph的教育智能体原型架构设计，实现核心教学状态机逻辑并接入大语言模型；完成RAG检索模块开发与初中物理知识库初步构建。 \\
        \hline
        2026年4月 & 集成防答案泄露护栏机制，开展系统内部调试；开发并配置模拟学生测试环境，形成测试样例集，开展初步功能测试与对抗性测试。 \\
        \hline
        2026年5月 & 开展仿真对话实验、基线对照与模块消融测试，收集对话日志并进行质量评估；根据实验结果优化提示策略、状态转移逻辑与护栏规则，完成专家评估准备。 \\
        \hline
        2026年6月 & 整理实验数据，完成结果分析与论文主体撰写，形成硕士学位论文初稿，开展后续修改与答辩准备工作。 \\
        \hline
    \end{tabular}
\end{table}


\noindent \textbf{预期成果：}
\begin{enumerate}
    \item 系统原型：可运行的、面向初中物理典型迷思概念（以电学与浮力为例）的苏格拉底式对话教育智能体原型。
    \item 知识与策略资源库：结构化的初中物理典型迷思概念条目库及其对应的苏格拉底式对话教学策略库。
    \item 学术论文：硕士学位论文《面向初中物理典型迷思概念（以电学与浮力为例）的苏格拉底式对话教育智能体设计与实现》。
\end{enumerate}


